\section{Trading Games}

\begin{frame}{What Are Trading Games?}
  \textbf{Setup}
  \begin{itemize}
    \item Players trade a stock whose true value is unknown
    \item Everyone starts with equal cash and position
    \item Players send bids and asks, which get matched by the engine
  \end{itemize}

  \vspace{0.3cm}
  \textbf{Information}
  \begin{itemize}
    \item Partial information about the true value is revealed over rounds
    \item Players update their estimates and trade accordingly
  \end{itemize}

  \vspace{0.3cm}
  \textbf{Scoring}
  \begin{itemize}
    \item At the end, true value is revealed, and all positions are liquidated at this price.
    \item Final Cash Position = cash + (position $\times$ true value)
  \end{itemize}
\end{frame}

\begin{frame}{Trading Game Example}
  \textbf{``The stock value is the volume of the tallest building in the world.''}

  \vspace{0.3cm}
  \textbf{Round 1}: ``The height of this building is 828 meters.''
  \begin{itemize}
    \item Players estimate volume and trade
  \end{itemize}

  \vspace{0.3cm}
  \textbf{Round 2}: ``The base area is 8,000 m$^2$.''
  \begin{itemize}
    \item Players refine estimates and trade
  \end{itemize}

  \vspace{0.3cm}
  \textbf{End}: True volume revealed, positions liquidated.
\end{frame}

\begin{frame}{Trading Game Formats}
  \textbf{How We Play}
  \begin{itemize}
    \item \textbf{Open outcry}: verbal, shouting quotes, manually tracking position on paper
    \item \textbf{Electronic}: enter orders into a simulated exchange
  \end{itemize}

  \vspace{0.3cm}
  \textbf{Why They Matter}
  \begin{itemize}
    \item Build intuition for markets, orderbooks, and spreads
    \item Improve estimation \& mental math skills
    \item Appears in interviews at some trading firms (Jane Street, Optiver)
  \end{itemize}
\end{frame}

\begin{frame}{Example Markets}
  \textbf{Make a market on...}
  \begin{itemize}
    \item Number of sheep in New Zealand
    \item Distance from Sydney to Perth in inches
    \item Number of 7-Elevens in Tokyo
  \end{itemize}

  \vspace{0.3cm}
  Extra information can be released as the market develops: 
  \begin{itemize}
    \item Catalysts (signals for change) 
    \item Noise (unrelated, no signal)
    \item Misinformation (misleading signals) 
  \end{itemize}

  \vspace{0.3cm}
\end{frame}

\begin{frame}{Open Outcry Rules}
  There are many variations of how you can play trading games without an electronic system.

  Any player can shout a quote at any time:
  \begin{itemize}
    \item \textbf{Bid [bid]} (i.e. insert bid at [bid])
    \item \textbf{Offer [offer]} (i.e. insert offer at [offer])
    \item \textbf{``[bid] at [offer]''} (i.e. insert both bid and offer)
  \end{itemize}

  \vspace{0.3cm}
  Trading against passive orders:
  \begin{itemize}
    \item \textbf{``Yours''}: trade against the highest bid (the asset is now ``yours'')
    \item \textbf{``Mine''}: trade against the lowest offer (the asset is now ``mine'')
  \end{itemize}
  \textcolor{sqt-question}{Keep track of your own PnL!}
\end{frame}

\begin{frame}{Electronic Rules}
  How to join:
  \begin{enumerate}
    \item Go to: \textcolor{sqt-question}{https://sqttrading.up.railway.app/register}
    \item Register using any name and password. You do not need to verify.
    \item Go to: \textcolor{sqt-question}{https://sqttrading.up.railway.app/compete}
    \item Review the interface and get ready to trade!
  \end{enumerate}
  Do not use any online resources or calculators.
\end{frame}
